%% ma: L12-B16-0023
%% lop: 12
%% chuong: 5
%% bai: 16
%% dang: 12.16.16
%% loai: DS
%% muc: [TH, TH, VD, VD]
%% dap_an: "SSSS"
%% nguon: "Ảnh giáo viên gửi 07/11/2026 (Tư Duy Mở - tuduymo.com, đề tự luyện Phương trình mặt phẳng, đề lập phương trình mặt phẳng), Câu 16"
%% kien_thuc: "Bài 14 (giao điểm với trục, phương trình đoạn chắn), Bài 16 (góc giữa hai mặt phẳng), thể tích tứ diện (lớp 11)"
%% trang_thai: chinh-thuc
%% ghi_chu: "Tài liệu gốc không kèm đáp án; đáp án do Claude giải và kiểm bằng sympy. Ý d dùng góc giữa hai mặt phẳng nên xếp vào Bài 16. Dạng đã duyệt 07/11/2026."
\begin{ex}
Trong không gian $Oxyz$, cho mặt phẳng $(P)\colon x+2y+3z-18=0$. Gọi $A$, $B$, $C$ lần lượt là giao điểm của $(P)$ với các trục tọa độ $Ox$, $Oy$, $Oz$.
\choiceTF
{Tọa độ của điểm $B$ là $(0;6;0)$.}
{Thể tích của tứ diện $OABC$ bằng $144$ (đvtt).}
{Có tất cả $20$ điểm có tọa độ nguyên dương thuộc mặt phẳng $(P)$.}
{Mặt phẳng $(\alpha)$ cắt ba tia $Ox$, $Oy$, $Oz$ lần lượt tại các điểm $M$, $N$, $P$ sao cho $O.MNP$ là hình chóp đều. Côsin của góc giữa $(\alpha)$ và $(P)$ bằng $\dfrac12$.}
\loigiai{$A(18;0;0)$, $B(0;9;0)$, $C(0;0;6)$.

a) $B(0;9;0)$. Sai.

b) $V=\dfrac16\cdot18\cdot9\cdot6=162$. Sai.

c) $x+2y+3z=18$ với $x,y,z\ge1$ nguyên: $z=1$ có $6$ điểm, $z=2$ có $5$, $z=3$ có $4$, $z=4$ có $3$, $z=5$ có $1$; tổng $19$ điểm. Sai.

d) $OM=ON=OP$ nên $(\alpha)\colon x+y+z=m$ ($m>0$). $\cos=\dfrac{|1+2+3|}{\sqrt3\cdot\sqrt{14}}=\dfrac{\sqrt{42}}7\ne\dfrac12$. Sai.}
\end{ex}
